% Inactive continuous Hessian foundations. No arbitrary bounded-product extension.
% Inputs: classical Garding algebra and smooth linear elliptic potential theory.
\section{Continuous Hessian products and domination}

Let $(X,\omega)$ be a connected compact Hermitian manifold of dimension
$n$, with $2\leq m<n$. Use the continuous admissibility condition of
604. Its linear tests imply $L_\omega u\geq-n$, where
$L_\omega u=n\ddc u\wedge\omega^{n-1}/\omega^n$.
We use the standard smooth elliptic maximum principle and Green estimate
\[
 \int_X(\sup_X u-u)\omega^n\leq C_\omega.
\]
The latter follows from the Green kernel for $L_\omega$, its positive
stationary probability, and $L_\omega u+n\geq0$. The stationary measure has
a smooth positive density equivalent to $\omega^n$, and the positive
measure $L_\omega u+n$ has fixed stationary mass $n$.
These are linear elliptic inputs, as in Ko\l odziej--Nguyen Lemma 3.3.
Maxima and convex combinations preserve continuous admissibility:
for each smooth positive linear test, a local smooth solution of the
inhomogeneous linear equation reduces this assertion to the maximum
property of subharmonic functions for that test operator.

For the algebra below, G\aa rding polarized positivity and convexity
are standard inputs. We also use the deleted-coordinate positivity
$S_{j;I}(\lambda)>0$ for $j+|I|\leq m$, a consequence of the same
polarized positivity. Indices of negative powers in derivative expansions
are omitted throughout.

\begin{lemma}\label{draft:605-absolute-mixed-forms}
If $0\leq r<m$, $\alpha_1,\ldots,\alpha_r\in\Gamma_m(\omega)$,
and $\Xi$ is a smooth complementary form, then
\[
 |\alpha_1\wedge\cdots\wedge\alpha_r\wedge\Xi|
 \leq C_{n,m,\Xi}(\alpha_1+\cdots+\alpha_r)^r\wedge\omega^{n-r}.
\]
For $r=0$ the right-hand side means $C_{n,m,\Xi}\omega^n$.
\end{lemma}
\begin{proof}
First take a single form $\alpha$, diagonal in an $\omega$-unitary frame,
with ordered eigenvalues $\lambda\in\Gamma_m$.
For $1\leq r<m$, repeated splitting and deleted-coordinate positivity
give $S_r(\lambda)\geq\lambda_1\cdots\lambda_r$.
The sum of the last $n-r$ eigenvalues is positive, so each negative
eigenvalue has modulus at most $(n-r)\lambda_{r+1}$.
Thus every product of $r$ distinct eigenvalues has modulus at most
\[
 (n-r)^r\lambda_1\cdots\lambda_r\leq(n-r)^r S_r(\lambda).
\]
Indeed, replace its negative factors by this bound and its positive
factors by the corresponding largest eigenvalues. In the diagonal frame
only complementary diagonal coefficients of $\Xi$ contribute to
$\alpha^r\wedge\Xi$. This proves
$|\alpha^r\wedge\Xi|\leq C_{n,m,\Xi}\alpha^r\wedge\omega^{n-r}$.

Apply this bound to the polynomial
$(\alpha_1+t_2\alpha_2+\cdots+t_r\alpha_r)^r\wedge\Xi$,
with $0\leq t_i\leq1$. Its absolute value is bounded by the stated
right-hand side, by polarized positivity and monotonicity on the cone.
A fixed invertible interpolation matrix on this finite-dimensional
polynomial space bounds each coefficient by its supremum on the cube.
The coefficient of $t_2\cdots t_r$ is a fixed positive multiple of
$\alpha_1\wedge\cdots\wedge\alpha_r\wedge\Xi$.
This proves the mixed bound. The closed-cone case follows by adding
$\varepsilon\omega$ to each form and taking the limit.
\end{proof}

\begin{proposition}\label{draft:605-continuous-products}
For continuous admissible $u_1,\ldots,u_k$, $1\leq k\leq m$, the current
\[
 (\omega+\ddc u_1)\wedge\cdots\wedge(\omega+\ddc u_k)
                                          \wedge\omega^{n-m}
\]
is well defined as the limit of uniform smooth admissible approximations.
It is positive, local, and continuous when its potentials converge
uniformly. Its local masses are bounded in terms of uniform bounds on
the potentials. In particular $H_m(u)$ is a positive Radon measure.
\end{proposition}
\begin{proof}
The first two lemmas of 604 provide the smooth approximations without
using any continuous Hessian product or domination input.
For a smooth admissible $u$ put $\sigma=\omega+\ddc u$.
Since $\mathrm d\sigma=\mathrm d\omega$, the coefficients of
$\ddc(\chi\sigma^{r-1}\wedge\omega^{n-r})$ are sums of smooth test
factors paired with powers $\sigma^q$, $q\leq r-1$, where $\chi$ is a
fixed cutoff. The absolute-form lemma and integration by parts give
the Chern--Levine--Nirenberg bound by induction on $r\leq m$:
\[
 \int_K\sigma^r\wedge\omega^{n-r}
                     \leq C_{K,U,\omega}(1+\|u\|_{L^\infty(U)})^r.
\]
At the induction step write $\sigma^r=\omega\wedge\sigma^{r-1}
+\ddc u\wedge\sigma^{r-1}$, integrate the second term against $\chi$,
and move $\ddc$ to $\chi\sigma^{r-1}\omega^{n-r}$.
The resulting powers have degrees $r-1,r-2,r-3$, and the induction
applies on nested slightly larger compact sets. Mixed masses reduce to
the single potential average, by positivity of the expansion of its
power. Positive mixed currents have total variation controlled by these
trace masses.

To prove independence, telescope the difference between two smooth
mixed products whose potentials have a common uniform bound. Each term
paired with a fixed smooth complementary test is
\[
 \int\ddc(u_a-v_a)\wedge T\wedge\eta
                         =\int(u_a-v_a)\ddc(T\wedge\eta),
\]
where $T$ is a product of the remaining admissible factors and
$\omega^{n-m}$. The identity $\mathrm d(\omega+\ddc u)=\mathrm d\omega$
and its $\dc$ analogue show that derivatives of $T$ introduce only
fixed metric derivatives; every term has at most $k-1<m$ admissible
factors. The absolute mixed-form bound and the mass estimate just
proved bound the variation of $\ddc(T\wedge\eta)$ uniformly.
Consequently the pairing difference is at most
$C_\eta\max_a\|u_a-v_a\|_\infty$.
The same construction with an arbitrary complementary smooth test
also defines the signed products of fewer than $m$ factors without the
factor $\omega^{n-m}$: their variation is bounded by the absolute-form
lemma. These products justify the metric-derivative terms above.
This proves that all approximations have the same limit and that the
products are continuous in the uniform topology. Positivity passes
from the smooth mixed products. Using a test supported in an open set
on which two limiting potentials coincide proves locality by the same
estimate on a smaller neighborhood of its support.
\end{proof}

\begin{lemma}\label{draft:605-contact-comparison}
If $u\leq v$ are continuous admissible functions and the remaining
$m-1$ factors are fixed continuous admissible functions, then on
$\{u=v\}$ the mixed measure with first factor $\omega+\ddc u$ is at
most the one with first factor $\omega+\ddc v$.
In particular $\mathbf1_{\{u=v\}}H_m(u)\leq
\mathbf1_{\{u=v\}}H_m(v)$.
\end{lemma}
\begin{proof}
Write $M(w)$ for the mixed measure with first potential $w$.
Positivity and locality give, for $t>0$,
\[
 M(\max(v,u+t))\geq
 \mathbf1_{\{v>u+t\}}M(v)+\mathbf1_{\{u\geq v\}}M(u).
\]
The first indicators increase to $\mathbf1_{\{v>u\}}$ as $t\downarrow0$,
and the potentials on the left converge uniformly to $v$.
Integrate against nonnegative continuous tests and pass to the limit.
Subtract the first restricted measure to obtain the claimed contact
inequality. Replacing one factor at a time proves the pure-power claim
and the analogous contact inequalities for all mixed powers.
\end{proof}

Define $P(h)$ to be the supremum of continuous admissible functions
bounded above by a continuous real obstacle $h$. Constants make this
class nonempty.

\begin{lemma}\label{draft:605-continuous-obstacle-envelope}
The function $P(h)$ is continuous and admissible, and $H_m(P(h))$ is
supported on $\{P(h)=h\}$. Moreover
$\|P(h_1)-P(h_2)\|_\infty\leq\|h_1-h_2\|_\infty$.
\end{lemma}
\begin{proof}
For smooth $h$, use the penalty solutions $w_N$ in the approximation
lemma 604. Its maximum argument applies to every continuous admissible
competitor $v\leq h$, with an error depending on $h$ and $N$, and not
on $v$. Taking their supremum gives
\[
 P(h)-E_N-2M_N/N\leq w_N\leq P(h),
\]
where the error tends to zero. Thus $P(h)$ is their uniform limit and
is continuous and admissible. On every compact subset of
$\{P(h)<h\}$ the penalty densities
$b e^{N(w_N-h)}$ tend uniformly to zero. Product continuity implies
the contact support assertion.
Adding constants to competitors gives the stated Lipschitz bound.
Approximate a continuous $h$ uniformly by smooth obstacles. Their
envelopes converge uniformly to $P(h)$, proving continuity and
admissibility. Any compact subset of $\{P(h)<h\}$ also lies outside the
contact sets of all sufficiently close smooth obstacles; product
continuity gives the support assertion for $h$.
\end{proof}

\begin{lemma}\label{draft:605-hessian-nonvanishing}
Every continuous admissible $u$ satisfies $\int_X H_m(u)>0$.
\end{lemma}
\begin{proof}
Suppose this integral is zero. For $s>0$ put
$w=\max(u,\min_Xu+s)$ and $U=\{u<\min_Xu+s\}$.
The open set $U$ is nonempty, $0\leq w-u\leq s$, and $w=u$ outside it.
Set $\sigma_u=\omega+\ddc u$, $\sigma_w=\omega+\ddc w$, and
$\mu_q=\sigma_u^q\wedge\sigma_w^{m-q}\wedge\omega^{n-m}$.
The contact lemma gives $\mathbf1_{X\setminus U}\mu_q\geq
\mathbf1_{X\setminus U}\mu_m$, for $0\leq q\leq m$.
For $q<m$ the telescoping identity is
\[
 \mu_q-\mu_m=\ddc(w-u)\wedge T_q,\qquad
 T_q=\sigma_u^q\wedge
       \sum_{r=0}^{m-q-1}\sigma_w^r\wedge\sigma_u^{m-q-1-r}
                                      \wedge\omega^{n-m}.
\]
Integration by parts on $X$ therefore gives
\[
 \int_U(\mu_q-\mu_m)\leq\int_X(\mu_q-\mu_m)
                                =\int_X(w-u)\ddc T_q.
\]
This identity follows first for smooth approximants and then by uniform
product continuity. The derivatives of $T_q$ contain only fixed metric
derivatives and products of at most $m-1$ admissible factors, so the
absolute-form lemma makes $\ddc T_q$ a measure of uniformly controlled
variation. Uniform convergence also justifies its continuous weight.
On $U$, locality gives $\sigma_w=\omega$ and thus
\[
 |\ddc T_q|\leq C\sum_{r=0}^{m-1}\mu_r.
\]
The constant depends only on the metric and $m,n$, and not on $u$ or $s$.
Consequently
\[
 \sum_{q=0}^m\int_U\mu_q
 \leq(m+1)\int_U H_m(u)+Cms\sum_{r=0}^{m-1}\int_U\mu_r.
\]
For sufficiently small $s$, absorption and $H_m(u)=0$ force all these
nonnegative integrals to vanish. But $\mu_0=\omega^n$ on the nonempty
open set $U$, a contradiction.
\end{proof}

\begin{proposition}\label{draft:605-zero-hessian-domination}
If continuous admissible $u,v$ satisfy
$\mathbf1_{\{u<v\}}H_m(u)=0$, then $u\geq v$.
\end{proposition}
\begin{proof}
Replace $v$ by $\max(u,v)$, so $u\leq v$.
For $b>1$ put $p_b=P(bu-(b-1)v)$ and
$z_b=b^{-1}p_b+(1-b^{-1})v\leq u$.
On the contact set $C_b=\{p_b=bu-(b-1)v\}$ one has $z_b=u$;
the contact inequality and polarized positivity give
\[
 \mathbf1_{C_b}b^{-m}H_m(p_b)
                 \leq\mathbf1_{C_b}H_m(z_b)
                 \leq\mathbf1_{C_b}H_m(u).
\]
The envelope measure is supported on $C_b$ and has positive mass.
It vanishes on $C_b\cap\{u<v\}$ by the hypothesis, so some point of
$C_b\cap\{u=v\}$ exists. Hence
$\sup_Xp_b\geq\inf_Xv$, while $p_b\leq u$.
The normalized linear $L^1$ estimate at the beginning is consequently
uniform for $p_b$. On the open set $E_a=\{u<v-a\}$, for $a>0$, one has
$p_b\leq v-ba$, and therefore
\[
 C_\omega\geq\int_X(\sup_Xp_b-p_b)\omega^n
                     \geq(ba-\operatorname{osc}_Xv)\int_{E_a}\omega^n.
\]
Letting $b\to\infty$ gives zero volume for $E_a$. Since it is open it
is empty, and $a\downarrow0$ gives $u\geq v$.
\end{proof}

\begin{proposition}\label{draft:605-strict-hessian-domination}
If continuous admissible $u,v$ satisfy
\[
 \mathbf1_{\{u<v\}}H_m(u)\leq
                    c\mathbf1_{\{u<v\}}H_m(v),\qquad 0\leq c<1,
\]
then $u\geq v$.
\end{proposition}
\begin{proof}
Choose $b>1$ so that $(1-b^{-1})^m>c$, and put
$p=P(bu-(b-1)v)$ and $z=b^{-1}p+(1-b^{-1})v\leq u$.
On $C=\{p=bu-(b-1)v\}$ the contact comparison gives
\[
 \mathbf1_C\bigl(b^{-m}H_m(p)+(1-b^{-1})^mH_m(v)\bigr)
                                  \leq\mathbf1_C H_m(u).
\]
Restricting to $\{u<v\}$ and using the hypothesis shows that
$H_m(p)$ vanishes on $C\cap\{u<v\}$. Its support lies in $C$, where
$p<v$ is equivalent to $u<v$. Thus $H_m(p)$ vanishes on $\{p<v\}$.
The preceding domination proposition gives $p\geq v$.
Since $p\leq bu-(b-1)v$, it follows that $u\geq v$.
\end{proof}

These continuous product and domination proofs supply the declared
pluripotential inputs in 604. They use the auxiliary existence and
approximation lemmas there, whose proofs depend only on the preceding
smooth estimates and elliptic theory. The stability and regularity
arguments follow afterwards, so no product or domination assertion is
used in its own construction. General bounded Hessian measures and
bounded-solution continuity remain outside this fragment.
